% ============================================================
% Round 03: PPO / RLHF
% ============================================================

\subsection{Round 03：PPO / RLHF}

\subsubsection{本轮问题}

\begin{conceptbox}
\textbf{学习目标：} 理解 Reward Model 的分数如何真正用来更新语言模型，以及 PPO 为什么需要 old policy、ratio、clip、value model 和 KL 约束。\\
\textbf{主线位置：}
\[
\text{SFT} \to \text{Reward Model} \to \text{PPO / RLHF} \to \text{DPO} \to \text{GRPO / RLVR}
\]
\end{conceptbox}

\subsubsection{从 RM 到 RLHF}

Reward Model 只会打分：

\[
R = r_\phi(x, y)
\]

但它不会直接改变语言模型。RLHF 要解决的问题是：

\[
\text{如何让 policy 更倾向于生成高 reward 的回答？}
\]

这里的 policy 就是语言模型：

\[
\pi_\theta(y \mid x)
\]

它给定 prompt $x$，生成回答 $y$。RLHF 的核心目标可以先粗略写成：

\[
\max_\theta \mathbb{E}_{y \sim \pi_\theta(\cdot \mid x)}
\left[
r_\phi(x, y)
\right]
\]

也就是：让模型自己采样回答，然后用 Reward Model 打分，再更新模型，使高分回答更容易被生成。

\begin{summarybox}
\textbf{一句话：} RLHF = 用 Reward Model 作为奖励函数，对 SFT 模型继续做强化学习。
\end{summarybox}

\subsubsection{把 LLM 看成 RL 问题}

在 LLM 里，一次生成可以看成一条 trajectory：

\[
y = (y_1, y_2, \ldots, y_T)
\]

对应 RL 里的状态、动作和奖励：

\begin{center}
{\small
\begin{tabular}{p{3cm}p{9.8cm}}
\hline
\textbf{RL 概念} & \textbf{LLM 中的含义} \\
\hline
State $s_t$ & prompt $x$ 加上已经生成的 tokens $y_{<t}$ \\
Action $a_t$ & 当前要生成的下一个 token $y_t$ \\
Policy $\pi_\theta$ & 语言模型的 next-token distribution \\
Trajectory & 一整段回答 $y_1,\ldots,y_T$ \\
Reward & Reward Model 对最终回答的评分，通常在序列末尾给出 \\
\hline
\end{tabular}
}
\end{center}

这说明 LLM 的 RL 和游戏 RL 有一个重要差别：奖励通常不是每一步都有，而是整段回答生成完以后才得到一个最终分数。

\subsubsection{为什么不能直接最大化 Reward}

如果只优化：

\[
\max_\theta \mathbb{E}[r_\phi(x, y)]
\]

模型可能会学会利用 Reward Model 的漏洞，而不是变得真正更好。例如：

\begin{itemize}[leftmargin=1.5em]
  \item 生成 Reward Model 偏爱的固定模板；
  \item 输出看起来很自信但事实错误的回答；
  \item 变得过度冗长，因为 RM 误以为长答案更有帮助；
  \item 偏离 SFT 模型的正常语言分布，导致可读性下降。
\end{itemize}

因此 RLHF 通常加入 reference model 约束。reference model 通常是冻结的 SFT 模型：

\[
\pi_\text{ref}
\]

实际优化目标更接近：

\[
\max_\theta
\mathbb{E}
\left[
r_\phi(x, y)
-
\beta
\mathbb{D}_\text{KL}
\left(
\pi_\theta(\cdot \mid x)
\|
\pi_\text{ref}(\cdot \mid x)
\right)
\right]
\]

\begin{intubox}
\textbf{直觉：} Reward Model 负责告诉模型“往哪里变好”，KL / reference model 负责告诉模型“不要跑太远”。
\end{intubox}

\subsubsection{PPO 的核心对象：概率比值}

PPO 更新时会区分两个 policy：

\begin{itemize}[leftmargin=1.5em]
  \item $\pi_{\theta_\text{old}}$：采样数据时使用的旧模型；
  \item $\pi_\theta$：当前正在更新的新模型。
\end{itemize}

对于每个 token，定义概率比值：

\[
\rho_t(\theta)
=
\frac{
\pi_\theta(y_t \mid x, y_{<t})
}{
\pi_{\theta_\text{old}}(y_t \mid x, y_{<t})
}
\]

它表示：新模型相比旧模型，是否更想生成这个 token。

\begin{center}
{\small
\begin{tabular}{p{3cm}p{9.8cm}}
\hline
\textbf{比值} & \textbf{含义} \\
\hline
$\rho_t > 1$ & 新模型提高了这个 token 的概率 \\
$\rho_t < 1$ & 新模型降低了这个 token 的概率 \\
$\rho_t \approx 1$ & 新旧模型变化不大 \\
\hline
\end{tabular}
}
\end{center}

\subsubsection{追问：new logprob 和 old logprob 分别是什么}

实际代码里通常不会直接除概率，而是先算 log probability：

\[
\rho_t(\theta)
=
\exp\left(
\log \pi_\theta(y_t \mid x, y_{<t})
-
\log \pi_{\theta_\text{old}}(y_t \mid x, y_{<t})
\right)
\]

也就是：

\[
\text{ratio}
=
\exp(\text{new\_logprob} - \text{old\_logprob})
\]

这里的分子和分母要看清楚：

\begin{center}
{\small
\begin{tabular}{p{3.2cm}p{9.6cm}}
\hline
\textbf{对象} & \textbf{含义} \\
\hline
\code{new\_logprob}
&
当前正在更新的 policy $\pi_\theta$ 对旧 rollout 里同一个 token $y_t$ 的 log probability：
$\log \pi_\theta(y_t \mid x, y_{<t})$。
\\
\code{old\_logprob}
&
采样这条 rollout 时的旧 policy $\pi_{\theta_\text{old}}$ 对同一个 token $y_t$ 的 log probability：
$\log \pi_{\theta_\text{old}}(y_t \mid x, y_{<t})$。
\\
\hline
\end{tabular}
}
\end{center}

\begin{warnbox}
\textbf{关键点：} PPO 更新时，新模型不是重新生成一个 token 再比较，而是拿旧 policy 已经采样出来的同一个 token，分别问：
\[
\text{旧模型当时给它多大概率？}
\qquad
\text{新模型现在给它多大概率？}
\]
两者相除，才得到 importance sampling ratio。
\end{warnbox}

所以一次 PPO rollout / update 可以理解成：

\begin{enumerate}[leftmargin=1.5em]
  \item 用 old policy $\pi_{\theta_\text{old}}$ 生成回答 $y_1,\ldots,y_T$；
  \item 在生成时保存每个 token 的 \code{old\_logprob}；
  \item PPO 更新时，把同一批 token 喂给当前 policy $\pi_\theta$；
  \item 重新计算这些 token 的 \code{new\_logprob}；
  \item 用 $\exp(\text{new\_logprob}-\text{old\_logprob})$ 得到 ratio。
\end{enumerate}

\begin{summarybox}
\textbf{不要混淆：} PPO ratio 的分母是 old policy，不是 reference model。\\
PPO 中常见的三个模型是：
\[
\pi_\theta
\quad
\pi_{\theta_\text{old}}
\quad
\pi_\text{ref}
\]
其中：
\begin{itemize}[leftmargin=1.5em]
  \item $\pi_\theta / \pi_{\theta_\text{old}}$ 用来计算 PPO importance sampling ratio；
  \item $\pi_\theta$ 和 $\pi_\text{ref}$ 的差异用来计算 KL penalty；
  \item $\pi_\text{ref}$ 通常是冻结的 SFT 模型，负责约束 policy 不要偏离太远。
\end{itemize}
\end{summarybox}

\subsubsection{追问：old policy 到底是什么模型}

old policy 不是一种新的模型结构。它就是\key{采样 rollout 时那一刻的 policy 快照}。

假设当前正在训练的模型参数是 $\theta$。在一次 PPO iteration 开始时，先把当前 policy 固定下来：

\[
\pi_{\theta_\text{old}} \leftarrow \pi_\theta
\]

然后用这个固定的 $\pi_{\theta_\text{old}}$ 去生成回答：

\[
y \sim \pi_{\theta_\text{old}}(\cdot \mid x)
\]

生成时保存每个 token 的 old logprob：

\[
\log \pi_{\theta_\text{old}}(y_t \mid x, y_{<t})
\]

接下来 PPO 更新的是新的 $\pi_\theta$。在这轮更新过程中，$\pi_{\theta_\text{old}}$ 不再变化，只作为分母提供基准：

\[
\rho_t(\theta)
=
\frac{
\pi_\theta(y_t \mid x, y_{<t})
}{
\pi_{\theta_\text{old}}(y_t \mid x, y_{<t})
}
\]

\begin{intubox}
\textbf{直觉：} old policy 就是“生成这批训练样本的那个旧版本模型”。PPO 问的是：当前新模型相比当时生成样本的旧模型，对同一个 token 的概率改了多少。
\end{intubox}

一次完整循环可以写成：

\begin{enumerate}[leftmargin=1.5em]
  \item 令 $\pi_{\theta_\text{old}} \leftarrow \pi_\theta$；
  \item 用 $\pi_{\theta_\text{old}}$ 采样一批回答，并保存 old logprob；
  \item 固定 $\pi_{\theta_\text{old}}$，用 PPO 多步更新 $\pi_\theta$；
  \item 这一批数据用完后，进入下一轮，再令新的当前 policy 变成下一轮的 old policy。
\end{enumerate}

\begin{warnbox}
\textbf{不要混淆三个角色：}
\begin{itemize}[leftmargin=1.5em]
  \item $\pi_\theta$：正在被训练、参数会更新的当前 policy；
  \item $\pi_{\theta_\text{old}}$：采样当前这批 rollout 时的 policy 快照，用于 PPO ratio 的分母；
  \item $\pi_\text{ref}$：冻结的 reference model，通常是 SFT 模型，用于 KL penalty。
\end{itemize}
第一轮 PPO 开始时，$\pi_\theta$ 通常从 SFT 模型初始化；但训练过程中，old policy 是每一轮采样前的当前 policy 快照，不等于永远固定不动的 reference model。
\end{warnbox}

\subsubsection{追问：如果没有多步更新，old policy 又是什么}

这里要区分两种说法。

\paragraph{情况一：每批 rollout 只做一次 PPO 更新。}

这是合理的。流程变成：

\begin{enumerate}[leftmargin=1.5em]
  \item 令 $\pi_{\theta_\text{old}} \leftarrow \pi_\theta$；
  \item 用 $\pi_{\theta_\text{old}}$ 采样一批回答；
  \item 用这批数据对 $\pi_\theta$ 做一次 PPO 更新；
  \item 丢掉这批 rollout；
  \item 下一轮重新令 $\pi_{\theta_\text{old}} \leftarrow \pi_\theta$，再采样新回答。
\end{enumerate}

这种情况下，old policy 仍然是\key{采样当前这批数据的 policy}。只不过这批数据只被用了一次，没有像常见 PPO 那样对同一批 rollout 做多个 epoch / minibatch 更新。

\begin{intubox}
\textbf{直觉：} 如果每次更新后都重新采样新数据，那么 old policy 每轮都会刷新。这仍然是 on-policy 的，只是样本利用率较低，训练成本更高。
\end{intubox}

\paragraph{情况二：还在用同一批 rollout，却每个梯度步都更新 old policy。}

这就不对。因为 PPO ratio 的分母必须是\key{生成这批 token 的那个 policy}：

\[
\rho_t(\theta)
=
\frac{
\pi_\theta(y_t \mid x, y_{<t})
}{
\pi_{\theta_\text{old}}(y_t \mid x, y_{<t})
}
\]

如果同一批 rollout 还没用完，就把 $\pi_{\theta_\text{old}}$ 也改成了最新的 $\pi_\theta$，那么分母不再代表“当时采样这些 token 的模型”。ratio 的重要性采样含义就被破坏了。

极端情况下，如果每个梯度步之后都令：

\[
\pi_{\theta_\text{old}} \leftarrow \pi_\theta
\]

并且继续用旧 rollout 训练，那么下一步的 ratio 会被人为拉回接近 $1$。这样 PPO-clip 就很难准确限制“当前 policy 相对采样 policy 走了多远”。

\begin{warnbox}
\textbf{核心判断：} old policy 可以在每一轮新采样前刷新；但不能在复用同一批 rollout 的中途随便刷新。\\
只要这批 token 还是由某个旧 policy 生成的，ratio 的分母就应该保持为那个生成它们的 policy。
\end{warnbox}

\begin{summarybox}
\textbf{最小结论：}
\[
\text{old policy} = \text{behavior policy} = \text{生成当前这批 rollout 的 policy}
\]
如果每次更新都重新采样，old policy 就每轮更新一次，这是可以的。\\
如果不重新采样、还在用旧数据，却把 old policy 跟着训练过程一起更新，那么 PPO ratio 就失去了正确含义。
\end{summarybox}

\subsubsection{追问：为什么有些 GRPO 实现里 ratio 看起来恒等于 1}

在 open-r1 使用的 TRL \code{GRPOTrainer} 里，有一类实现会出现类似写法：

\[
\text{coef}
=
\exp\left(
\text{per\_token\_logps}
-
\text{per\_token\_logps.detach()}
\right)
\]

前向计算时：

\[
\text{per\_token\_logps}
-
\text{per\_token\_logps.detach()}
=
0
\]

所以：

\[
\text{coef} = \exp(0) = 1
\]

这就是“importance ratio 看起来恒等于 1”的来源。

但这个写法不能只看前向值。因为 \code{detach()} 会切断分母那一侧的梯度：

\[
\nabla_\theta
\left[
\exp\left(
\log \pi_\theta(y_t \mid s_t)
-
\operatorname{stopgrad}(\log \pi_\theta(y_t \mid s_t))
\right)
A_t
\right]
=
A_t \nabla_\theta \log \pi_\theta(y_t \mid s_t)
\]

也就是说：

\begin{itemize}[leftmargin=1.5em]
  \item 前向数值上，ratio 是 $1$；
  \item 反向传播时，它仍然给出 policy gradient；
  \item 这个 loss 的数值本身不重要，重要的是它的梯度。
\end{itemize}

\begin{intubox}
\textbf{直觉：} 这不是在做“新 policy 除以旧 policy”的完整 PPO 多步 ratio，而是在没有复用旧 rollout 多步更新时，用一个 stop-gradient 技巧写出 REINFORCE 风格的梯度：
\[
-A_t \nabla_\theta \log \pi_\theta(y_t \mid s_t)
\]
\end{intubox}

这通常对应一个特殊训练设置：

\begin{center}
{\small
\begin{tabular}{p{4cm}p{8.8cm}}
\hline
\textbf{训练设置} & \textbf{old logprob 的处理} \\
\hline
每批 rollout 只更新一次，不做多轮复用
&
old policy 和当前 policy 在这次更新开始时相同，可以用 \code{per\_token\_logps.detach()} 作为分母；前向 ratio 为 $1$，但仍有 policy gradient。
\\
同一批 rollout 要复用多次更新
&
必须保存采样时的 \code{old\_per\_token\_logps}，否则 ratio 无法衡量当前 policy 相对采样 policy 走了多远。
\\
\hline
\end{tabular}
}
\end{center}

\begin{warnbox}
\textbf{关键区别：} ratio 前向为 $1$ 不等于“没有训练”。\\
如果写成普通常数 $1 \cdot A_t$，确实没有 policy 梯度；但如果写成
\[
\exp(\log p - \operatorname{stopgrad}(\log p)) A_t
\]
前向值是 $A_t$，梯度却包含 $A_t \nabla \log p$。
\end{warnbox}

\begin{summarybox}
\textbf{最小结论：} open-r1 / TRL 里看到 ratio 为 $1$，一般说明它处在“没有多步复用旧 rollout”的简化场景，或者使用了 stop-gradient 技巧。\\
它和标准 PPO 里需要显式保存
\[
\log \pi_{\theta_\text{old}}(y_t \mid s_t)
\]
来做多步 clipped update 的场景不是一回事。
\end{summarybox}

\subsubsection{Advantage：这个回答该鼓励还是惩罚}

Reward 本身只是一个分数。PPO 需要知道某个 token / trajectory 相比预期到底好多少，这就是 Advantage：

\[
A_t = R_t - V_\psi(s_t)
\]

其中 $V_\psi(s_t)$ 是 value model / critic，估计当前状态后续大概能拿到多少 reward。

\begin{itemize}[leftmargin=1.5em]
  \item $A_t > 0$：比预期更好，应该提高概率；
  \item $A_t < 0$：比预期更差，应该降低概率；
  \item $A_t \approx 0$：和预期差不多，更新很小。
\end{itemize}

\begin{warnbox}
\textbf{注意：} PPO 需要额外训练 value model / critic。这会增加显存、训练复杂度和不稳定性。后面的 GRPO 一个核心动机就是去掉 critic，用组内相对奖励估计 advantage。
\end{warnbox}

\subsubsection{PPO-Clip 目标函数}

PPO 的基本思想是：如果某个 token 来自高 advantage 的回答，就提高它的概率；如果来自低 advantage 的回答，就降低它的概率。但每次更新不能太大。

clip 目标函数：

\[
\mathcal{J}_\text{PPO}(\theta)
=
\mathbb{E}_t
\left[
\min
\left(
\rho_t(\theta) A_t,\;
\operatorname{clip}(\rho_t(\theta), 1-\epsilon, 1+\epsilon) A_t
\right)
\right]
\]

其中 $\epsilon$ 通常是一个小常数，例如 $0.1$ 或 $0.2$。

\begin{intubox}
\textbf{直觉：} PPO 允许模型朝高 reward 方向移动，但用 clip 限制每一步别走太远。它解决的是 policy gradient 更新过猛、训练崩掉的问题。
\end{intubox}

\subsubsection{RLHF 的典型流程}

\begin{enumerate}[leftmargin=1.5em]
  \item 从 prompt dataset 采样一批 prompts；
  \item 用当前 policy $\pi_{\theta_\text{old}}$ 生成回答；
  \item 用 Reward Model 给回答打分；
  \item 计算 KL penalty，得到最终 reward；
  \item 用 value model 估计 advantage；
  \item 用 PPO-clip 更新 policy；
  \item 同时训练或更新 value model；
  \item 重复采样和更新。
\end{enumerate}

简化成一条链：

\[
\text{Prompt}
\to
\text{Policy 生成}
\to
\text{RM 打分}
\to
\text{KL 约束}
\to
\text{Advantage}
\to
\text{PPO 更新}
\]

\subsubsection{PPO / RLHF 的代价}

PPO/RLHF 很强，但工程代价高：

\begin{itemize}[leftmargin=1.5em]
  \item 需要 policy model；
  \item 需要 reference model；
  \item 需要 reward model；
  \item 通常还需要 value model / critic；
  \item 训练时要保存 log probability、reward、advantage 等中间量；
  \item 超参数敏感，例如 KL 系数、clip range、reward normalization。
\end{itemize}

这也是为什么后来出现 DPO 和 GRPO：

\begin{itemize}[leftmargin=1.5em]
  \item DPO：希望不用显式训练 RM、也不用在线 PPO；
  \item GRPO：希望保留在线 RL 的能力，但去掉 critic。
\end{itemize}

\subsubsection{本轮最小结论}

\begin{summarybox}
\textbf{PPO/RLHF = 用 RM 分数训练 policy，但用 clip 和 KL 控制更新幅度。}\\
关键对象：
\[
\rho_t(\theta)
=
\frac{
\pi_\theta(y_t \mid x, y_{<t})
}{
\pi_{\theta_\text{old}}(y_t \mid x, y_{<t})
}
\]
关键思想：高 advantage 的 token 增大概率，低 advantage 的 token 降低概率；clip 和 KL 防止模型为了 reward 跑偏。
\end{summarybox}

\subsubsection{本轮检查题}

\begin{enumerate}[leftmargin=1.5em]
  \item 在 LLM RLHF 里，state、action、trajectory、reward 分别是什么？
  \item PPO 里的 $\rho_t(\theta)$ 表示什么？为什么需要 old policy？
  \item Advantage $A_t$ 的作用是什么？为什么只看 reward 不够？
  \item PPO-clip 为什么能提高训练稳定性？
  \item KL / reference model 在 RLHF 里起什么作用？
  \item PPO/RLHF 的工程代价在哪里？这如何引出 DPO 和 GRPO？
\end{enumerate}
